%% GRB classification from compact binary mergers
%% AASTeX v7.0.2 manuscript, two-column ApJ style. Build:
%%   pdflatex grb_paper && bibtex grb_paper && pdflatex grb_paper && pdflatex grb_paper
%% Figures are read from ../plots/ via \graphicspath.

\documentclass[twocolumn,amsmath]{aastex702}

\graphicspath{{../plots/}{./}}

%% Running heads
\shorttitle{Compact binary mergers to GRB populations}
\shortauthors{Rodriguez, Levina, \& Broekgaarden}

%% Project macros
\newcommand{\Msun}{M_\odot}
\newcommand{\Mtov}{M_{\rm TOV}}
\newcommand{\Mthresh}{M_{\rm thresh}}
\newcommand{\Mtot}{M_{\rm tot}}
\newcommand{\Mdisk}{M_{\rm disk}}
\newcommand{\compas}{\textsc{compas}}
\newcommand{\fci}{\textsc{FastCosmicIntegration}}
\newcommand{\Rz}{\mathcal{R}(z)}
\newcommand{\Rzero}{\mathcal{R}(z{=}0)}
\newcommand{\aBH}{a_{\rm BH}}
\newcommand{\Gpcyr}{{\rm Gpc^{-3}\,yr^{-1}}}
\newcommand{\Zsun}{Z_\odot}

%% Float placement tuning: let pages hold more figures so full-width
%% figure* floats sit near their text instead of piling onto float pages.
%% Wide floats only ever land at the top of a page (or a float page), so
%% allow three per page and let figure-dominated pages through.
\setcounter{topnumber}{2}
\setcounter{dbltopnumber}{3}
\renewcommand{\topfraction}{0.9}
\renewcommand{\dbltopfraction}{0.95}
\renewcommand{\bottomfraction}{0.6}
\renewcommand{\textfraction}{0.07}
\renewcommand{\floatpagefraction}{0.7}
\renewcommand{\dblfloatpagefraction}{0.65}
\usepackage{placeins}

\begin{document}

\title{From Compact Binary Mergers to Gamma-Ray Bursts: Predicting GRB Populations Across Cosmic Time}

%% AASTeX v7 requires an \email for every \author.
\author[orcid=0009-0000-4068-2032]{Joseph Rodriguez}
\affiliation{Department of Astronomy \& Astrophysics, University of California, San Diego, 9500 Gilman Drive, La Jolla, CA 92093, USA}
\email[show]{jrodriguezruelas@ucsd.edu}

\author[orcid=0000-0003-1241-7615]{Sasha Levina}
\affiliation{Department of Astronomy \& Astrophysics, University of California, San Diego, 9500 Gilman Drive, La Jolla, CA 92093, USA}
\email{slevina@ucsd.edu}

\author[orcid=0000-0002-4421-4962]{Floor S. Broekgaarden}
\affiliation{Department of Astronomy \& Astrophysics, University of California, San Diego, 9500 Gilman Drive, La Jolla, CA 92093, USA}
\email{fbroekgaarden@ucsd.edu}

\begin{abstract}
Short and long gamma-ray bursts (GRBs) with kilonovae both trace compact binary
mergers; in the \citet{gottlieb2023,gottlieb2024} framework the difference is
set not by duration but by the central engine, a hyper-massive neutron star
(HMNS) or a black-hole accretion disk. We classify the \compas\
population-synthesis models of \citet{broekgaarden2021} into these engine
classes, propagating STROOPWAFEL importance weights through every reduction.
Intrinsic merger rates follow from cosmic integration over the
\citet{levina2026} IllustrisTNG TNG100-1 metallicity-dependent star-formation
history under Planck 2015 cosmology. Four results follow. First,
long-duration bursts dominate: 76 percent of binary neutron star (BNS) mergers
produce one (27 percent from an HMNS, 2 percent from a disk, 47 percent
faint), against a 24 percent short-duration minority with blue kilonovae. The
black hole neutron star (BHNS) population is instead dominated by silent
mergers: 76 percent swallow the neutron star and launch no burst at all. Second, the BNS local rate ($55\,\Gpcyr$) falls well inside the LVK
GWTC-5.0 90 percent credible region, while the BHNS rate ($84.5\,\Gpcyr$) sits
a factor $\sim 2.6$ above the tightened NSBH region; all five
population-synthesis variations exceed that region, so the \compas\ BHNS
population generically overpredicts the GW-measured NSBH rate. Third, against the \citet{rastinejad2024} kilonova sample,
the long-duration bursts match the black-hole-disk prediction to $0.1$ dex
while the short-duration bursts lie roughly three orders of magnitude above
it: the engine, not the duration, sets the kilonova. Fourth, across the full
20-model grid the local rate spans two orders of magnitude while the class
fractions barely move; the class demographics, not the absolute rates, are the
robust prediction. All GRB rates are upper bounds, assuming complete jet
launching above threshold.
\end{abstract}

\keywords{\uat{Gamma-ray bursts}{629}, \uat{Neutron stars}{1108}, \uat{Compact binary stars}{283}, \uat{Gravitational waves}{678}, \uat{R-process}{1324}, \uat{Stellar populations}{1622}}

%% ======================================================================
\section{Introduction} \label{sec:intro}
%% ======================================================================

Compact binary mergers that contain at least one neutron star (NS) power both
gravitational-wave (GW) transients and a rich set of electromagnetic
counterparts: short gamma-ray bursts (sGRBs), kilonovae, and, as the multi-band
sample has grown, a population of longer-duration bursts accompanied by red
kilonovae \citep{rastinejad2024,levan2024}. The joint detection of GW170817,
GRB\,170817A, and AT2017gfo established the binary neutron star (BNS) channel as
a confirmed sGRB progenitor \citep{abbott2019gw170817,goldstein2017,kasen2017},
and the second confirmed BNS GW190425 \citep{abbott2020gw190425} extended the
known mass range. A central open question is which central engine launches the
relativistic jet and sets the composition of the kilonova ejecta.

\citet{gottlieb2023,gottlieb2024} proposed that the merger remnant determines
the burst phenomenology along two physically distinct engine branches. A
black-hole (BH) accretion disk, formed either promptly in a BNS merger or after
NS tidal disruption in a black hole neutron star (BHNS) binary, launches a jet
whose energetics scale with the disk mass. A long-lived hyper-massive neutron
star (HMNS) or magnetar remnant instead drives a magnetically dominated wind
that both extends the burst duration and enriches the ejecta. The two branches
predict different durations, ejecta masses, and kilonova colors, and therefore
admit a population-level test once a synthetic compact-binary population is
classified and its cosmic rate computed.

In this paper we apply the Gottlieb classification to the \compas\
\citep{broekgaarden2021} population-synthesis models, centered on a five-member
core suite spanning the two systematics that most strongly affect the remnant
mass distribution, the common-envelope efficiency $\alpha_{\rm CE}$ and the
maximum neutron-star mass $M_{\rm NS,max}$, and extended to the full
20-variation grid in the appendices. We propagate the STROOPWAFEL importance
weights of the underlying simulation through every reduction, integrate the
intrinsic cosmic merger rate with the IllustrisTNG-calibrated metallicity-specific
star-formation history of \citet{levina2026}, and test the resulting local
rates against the LVK GWTC-5.0 credible regions \citep{abac2026}.

This paper makes four contributions: (i) per-class predictions of the BNS and
BHNS mass, delay-time, and metallicity distributions
(Sections~\ref{sec:massplane} to \ref{sec:delaytime}); (ii) intrinsic cosmic
merger rates resolved by GRB class and by \citet{broekgaarden2021} formation
channel, tested against the GWTC-5.0 credible regions
(Sections~\ref{sec:classrates} to \ref{sec:lvk}, with the channel
decomposition in Appendix~\ref{app:channels}); (iii) a confrontation of the
engine classes with the \citet{rastinejad2024} kilonova-bearing GRB sample,
through class fractions, through an engine discriminator built on burst
duration and kilonova ejecta mass, and through host-galaxy offset predictions
(Sections~\ref{sec:obsclass} to \ref{sec:offsets}); and (iv) a BH-spin
sensitivity analysis (Section~\ref{sec:spin}) plus a population-synthesis
robustness study spanning the five-model suite and the full 20-model
\citet{broekgaarden2021} grid (Appendices~\ref{app:robustness} and
\ref{app:grid}). Methods are collected in Section~\ref{sec:methods}, results
in Section~\ref{sec:results}, discussion in Section~\ref{sec:discussion}, and
conclusions in Section~\ref{sec:conclusions}.

%% ======================================================================
\section{Methods} \label{sec:methods}
%% ======================================================================

\subsection{COMPAS populations and importance weights} \label{sec:populations}

We use the \compas\ rapid binary population-synthesis catalogs of
\citet{broekgaarden2021,broekgaarden2022}, downloaded from Zenodo (BNS record
5189849, BHNS record 5178777). Both records carry the 20-variation paper II
data release of \citet{broekgaarden2022} and package each catalog under that
paper's letter convention; we label models by the paper I letters of
\citet{broekgaarden2021}, and Table~\ref{tab:models} maps one convention onto
the other. The release spans 20 binary-physics variations, each with a BNS
and a BHNS catalog (40 catalogs in total); Table~\ref{tab:models} collects
the assumption each variation changes relative to the fiducial, together
with the size of each catalog. The main text works with a five-member core
suite that isolates the two assumptions with the strongest leverage on the remnant
mass distribution: the fiducial Model A ($\alpha_{\rm CE} = 1.0$, maximum
neutron-star gravitational mass $M_{\rm NS,max} = 2.5\,\Msun$), F
($\alpha_{\rm CE} = 0.5$) and G ($\alpha_{\rm CE} = 2.0$) bracketing the
common-envelope efficiency, and J ($M_{\rm NS,max} = 2.0\,\Msun$) and K
($M_{\rm NS,max} = 3.0\,\Msun$) bracketing the maximum NS mass
(Appendix~\ref{app:robustness}); the remaining 15 variations enter the
grid-wide robustness scan of Appendix~\ref{app:grid}. All variations except
Model I (rapid) share the \citet{fryer2012} delayed supernova engine.
Common-envelope evolution follows the \citet{webbink1984} $\alpha$-formalism
with the \citet{xuli2010} binding-energy parameter $\lambda_{\rm CE}$ and
\citet{hurley2002} stellar tracks.

\begin{deluxetable*}{lclcclrr}
\tabletypesize{\footnotesize}
\tablecaption{The 20 binary-physics variations of the
\citet{broekgaarden2021,broekgaarden2022} grid. \label{tab:models}}
\tablehead{
\colhead{Model} & \colhead{Paper II} & \colhead{Varied assumption} & \colhead{$\alpha_{\rm CE}$} &
\colhead{$M_{\rm NS,max}$} & \colhead{SN remnant engine} &
\colhead{BNS $N$ ($n_{\rm eff}$)} & \colhead{BHNS $N$ ($n_{\rm eff}$)} \\
\colhead{} & \colhead{} & \colhead{} & \colhead{} & \colhead{($\Msun$)} & \colhead{} & \colhead{} & \colhead{}
}
\startdata
A\tablenotemark{a}        & A & Fiducial                                            & 1.0 & 2.5 & delayed & $233{,}137$ ($65{,}134$)      & $1{,}525{,}553$ ($292{,}484$) \\
alpha0p1\tablenotemark{b} & G & $\alpha_{\rm CE} = 0.1$                             & 0.1 & 2.5 & delayed & $345{,}779$ ($207{,}909$)     & $22{,}395$ ($7{,}418$) \\
F\tablenotemark{a}        & H & $\alpha_{\rm CE} = 0.5$                             & 0.5 & 2.5 & delayed & $64{,}260$ ($16{,}488$)       & $915{,}179$ ($113{,}902$) \\
G\tablenotemark{a}        & I & $\alpha_{\rm CE} = 2.0$                             & 2.0 & 2.5 & delayed & $1{,}052{,}472$ ($399{,}336$) & $833{,}433$ ($385{,}378$) \\
alpha10\tablenotemark{b}  & J & $\alpha_{\rm CE} = 10$                              & 10  & 2.5 & delayed & $579{,}059$ ($330{,}218$)     & $441{,}924$ ($278{,}983$) \\
B & B & Mass-transfer efficiency $\beta = 0.25$                                     & 1.0 & 2.5 & delayed & $41{,}984$ ($11{,}183$)       & $738{,}537$ ($150{,}896$) \\
C & C & Mass-transfer efficiency $\beta = 0.5$                                      & 1.0 & 2.5 & delayed & $45{,}070$ ($9{,}265$)        & $148{,}043$ ($63{,}576$) \\
D & D & Mass-transfer efficiency $\beta = 0.75$                                     & 1.0 & 2.5 & delayed & $292{,}349$ ($73{,}028$)      & $118{,}921$ ($45{,}300$) \\
E & E & Case-BB mass transfer unstable                                              & 1.0 & 2.5 & delayed & $371$ ($217$)                 & $458{,}667$ ($51{,}821$) \\
EH\tablenotemark{b} & F & Unstable case BB plus optimistic CE                       & 1.0 & 2.5 & delayed & $7{,}491$ ($4{,}490$)         & $828{,}558$ ($44{,}525$) \\
H & K & Optimistic CE (HG donors survive CE)                                        & 1.0 & 2.5 & delayed & $234{,}963$ ($35{,}646$)      & $1{,}535{,}042$ ($55{,}884$) \\
I & L & Rapid SN remnant prescription                                               & 1.0 & 2.5 & rapid   & $146{,}467$ ($43{,}297$)      & $2{,}766{,}298$ ($974{,}073$) \\
J\tablenotemark{a}        & M & $M_{\rm NS,max} = 2.0\,\Msun$                       & 1.0 & 2.0 & delayed & $213{,}669$ ($58{,}125$)      & $959{,}796$ ($189{,}149$) \\
K\tablenotemark{a}        & N & $M_{\rm NS,max} = 3.0\,\Msun$                       & 1.0 & 3.0 & delayed & $238{,}111$ ($68{,}058$)      & $1{,}990{,}330$ ($423{,}846$) \\
L & O & No (pulsational) pair-instability SNe                                       & 1.0 & 2.5 & delayed & $233{,}190$ ($65{,}787$)      & $1{,}524{,}497$ ($289{,}381$) \\
M & P & CCSN kick $\sigma_{\rm kick} = 100\,{\rm km\,s^{-1}}$                       & 1.0 & 2.5 & delayed & $415{,}335$ ($107{,}644$)     & $3{,}049{,}458$ ($691{,}439$) \\
N & Q & CCSN kick $\sigma_{\rm kick} = 30\,{\rm km\,s^{-1}}$                        & 1.0 & 2.5 & delayed & $626{,}019$ ($180{,}711$)     & $4{,}198{,}238$ ($1{,}077{,}213$) \\
O & R & No BH natal kick                                                            & 1.0 & 2.5 & delayed & $155{,}208$ ($51{,}465$)      & $5{,}068{,}628$ ($1{,}685{,}390$) \\
fWR0p1\tablenotemark{b}   & S & Wolf-Rayet wind multiplier $f_{\rm WR} = 0.1$       & 1.0 & 2.5 & delayed & $191{,}544$ ($58{,}259$)      & $1{,}379{,}487$ ($301{,}322$) \\
fWR5\tablenotemark{b}     & T & Wolf-Rayet wind multiplier $f_{\rm WR} = 5$         & 1.0 & 2.5 & delayed & $276{,}279$ ($71{,}624$)      & $938{,}755$ ($198{,}429$) \\
\enddata
\tablenotetext{a}{Core suite used throughout the main text and
Appendix~\ref{app:robustness}.}
\tablenotetext{b}{Variation introduced in the \citet{broekgaarden2022} paper II
release with no paper I letter; labeled by its descriptive token.}
\tablecomments{Rows are grouped by physics family (common envelope, mass
transfer, case-BB stability, CE survival, SN engine, NS mass cap, pair
instability, natal kicks, winds); Figure~\ref{fig:rate20} lists the core
suite first instead, and Figure~\ref{fig:frac20} sorts models by class
fraction. Each variation changes one
assumption relative to Model A; unvaried assumptions keep the fiducial values
($\alpha_{\rm CE} = 1.0$, adaptive $\beta$, stable case-BB mass transfer,
pessimistic CE survival, $M_{\rm NS,max} = 2.5\,\Msun$,
$\sigma_{\rm kick} = 265\,{\rm km\,s^{-1}}$, \citealt{hobbs2005},
$f_{\rm WR} = 1$). Each variation has both a BNS and a BHNS catalog,
40 catalogs in total, public on Zenodo (BNS record 5189849, BHNS record
5178777). The Model column uses the \citet{broekgaarden2021} paper I letters
(descriptive tokens for the five paper-II-only variations); the Paper II
column gives the \citet{broekgaarden2022} letter under which both Zenodo
records package the same variation. The two conventions diverge after model E
(paper II K is the optimistic-CE variation, not $M_{\rm NS,max} = 3.0\,\Msun$),
so files retrieved from the records must be matched by the Paper II letter.
$N$ is the raw number of merging systems in the catalog and $n_{\rm eff}$ the
Kish effective sample size of its STROOPWAFEL weights
(Section~\ref{sec:populations}); the Model A entries match
Table~\ref{tab:classfrac}.}
\end{deluxetable*}

The \compas\ samples are drawn under the STROOPWAFEL adaptive importance-sampling
algorithm \citep{broekgaarden2019}, so unweighted statistics are biased. Every reduction in this work uses the
per-system weights $w$: counts are $\sum_i w_i$, means are weighted means, and
histograms and kernel density estimates are weighted. Where a sample size enters
a figure caption we quote both the raw count $N$ and the Kish effective sample
size $n_{\rm eff} = (\sum_i w_i)^2 / \sum_i w_i^2$ \citep{kish1965}; the
per-catalog $N$ and $n_{\rm eff}$ of all 40 catalogs are listed in
Table~\ref{tab:models}. Masses are returned with the convention $m_1 \geq m_2$.

\subsection{Neutron-star mass remapping} \label{sec:remap}

Both \compas\ supernova engines inherit a $\sim 1.7\,\Msun$ deficit in the
neutron-star mass distribution from the \citet{fryer2012} baryonic-to-gravitational
mass conversion. Following \citet{mandel2020}, we close this gap by remapping the
synthetic NS masses onto the empirical double-Gaussian mass function of
\citet{alsing2018},
\begin{equation}
f(m) \propto W_1\,\mathcal{N}(m; \mu_1, \sigma_1) + W_2\,\mathcal{N}(m; \mu_2, \sigma_2),
\label{eq:nsremap}
\end{equation}
truncated to $[1.10,\ \Mtov]\,\Msun$, with $(W_1, \mu_1, \sigma_1) = (0.66, 1.34\,\Msun, 0.07\,\Msun)$
and $(W_2, \mu_2, \sigma_2) = (0.34, 1.80\,\Msun, 0.21\,\Msun)$. The first component is the
recycled and slow-pulsar peak near $1.34\,\Msun$ and the second is the high-mass
component near $1.80\,\Msun$; the values reproduce the \citet{alsing2018} Galactic
double-Gaussian fit. We map each (weighted) input mass through the empirical CDF onto
the quantile of this truncated target. For BNS we remap each component-mass
marginal in its own rank stream (stacking the two components into one stream
imprints an artificial deficit at the target median) and re-sort to
preserve $m_1 \geq m_2$; for BHNS we remap the marginal NS mass. The remap uses
fixed random seeds (42 for BNS, 43 for BHNS) for reproducibility, caps the NS mass
at $\Mtov$, and is applied once immediately after loading so that every downstream
classifier and rate uses the remapped masses. The remap shifts the total-mass
distribution and therefore changes class fractions relative to the raw catalog.

\subsection{Remnant and ejecta physics} \label{sec:remnant}

The BHNS classification rests on the post-merger accretion-disk mass, which we
assemble from a chain of numerical-relativity fitting formulae. The neutron-star
radius follows a phenomenological mass-radius relation anchored at
$R_{1.4} \approx 12$ km \citep{raaijmakers2021}, with a cubic softening toward $\Mtov$,
\begin{equation}
R_{\rm NS}(M) = R_{1.4}\left[1 - 0.15\left(\frac{M - 1.4\,\Msun}{\Mtov - 1.4\,\Msun}\right)^{3}\right],
\label{eq:nsradius}
\end{equation}
held flat at $R_{1.4}$ below $1.4\,\Msun$ and undefined above $\Mtov$. This is a
modeling choice rather than an equation-of-state derivation. The neutron-star
compactness is then
\begin{equation}
C = \frac{G M_{\rm NS}}{R_{\rm NS}\,c^2},
\label{eq:compactness}
\end{equation}
and the baryonic mass follows the quadratic approximation
\begin{equation}
M_b = M_{\rm NS} + 0.080\,M_{\rm NS}^2
\label{eq:baryon}
\end{equation}
\citep{gao2020,lattimer2001}, with masses in $\Msun$. For a Kerr black hole of dimensionless spin
$\chi \equiv \aBH$, the innermost stable circular orbit (ISCO) radius in units of
$G M_{\rm BH}/c^2$ is \citep{bardeen1972}
\begin{align}
Z_1 &= 1 + (1-\chi^2)^{1/3}\!\left[(1+\chi)^{1/3} + (1-\chi)^{1/3}\right], \nonumber\\
Z_2 &= \sqrt{3\chi^2 + Z_1^2}, \label{eq:isco}\\
\hat{r}_{\rm ISCO} &= 3 + Z_2 - \mathrm{sgn}(\chi)\sqrt{(3-Z_1)(3+Z_1+2Z_2)}, \nonumber
\end{align}
taking the prograde root for $\chi \geq 0$ and the retrograde root for $\chi < 0$.

The baryon mass left outside the black hole after tidal disruption is the
\citet[their Eq.~4]{foucart2018} fit
\begin{equation}
\frac{M_{\rm rem}}{M_b} = \left[\max\!\left(0,\ \alpha\,\frac{1-2C}{\eta^{1/3}}
- \beta\,\hat{r}_{\rm ISCO}\,\frac{C}{\eta} + \gamma\right)\right]^{\delta},
\label{eq:foucart}
\end{equation}
with symmetric mass ratio $\eta = M_{\rm NS} M_{\rm BH}/(M_{\rm NS}+M_{\rm BH})^2$,
mass ratio $Q \equiv M_{\rm BH}/M_{\rm NS}$, and best-fit coefficients
$(\alpha,\beta,\gamma,\delta) = (0.406, 0.139, 0.255, 1.761)$. The fit is
calibrated for $Q \in [1, 7]$, $\chi \in [-0.5, 0.9]$, and $C \in [0.13, 0.182]$.
We clip spins at $\chi = 0.9$ to stay inside the calibrated spin range; the
mass-ratio range has no such guard, and 63 percent of the merging BHNS sample
(43 percent by weight) has $Q > 7$, extending to $Q = 29.4$, where
Eq.~\ref{eq:foucart} is an extrapolation. The extrapolation is benign for the
classification because the fit predicts essentially no remnant mass at these
mass ratios. By weight, 57 percent of the No-GRB class sits at $Q > 7$,
against only 0.2 percent of the faint-lbGRB class and none of the lbGRB plus
red kilonova class; forcing every $Q > 7$ system into the No-GRB class would
change the class fractions by at most 0.05 percentage points.
The dynamical ejecta follow \citet[their Eq.~9]{kruger2020},
\begin{equation}
\frac{M_{\rm dyn}}{M_b} = \max\!\left(0,\ a_1 Q^{n_1}\frac{1-2C}{C}
- a_2 Q^{n_2}\hat{r}_{\rm ISCO} + a_4\right),
\label{eq:kfdyn}
\end{equation}
with $(a_1, a_2, a_4) = (0.007116, 0.001436, -0.02762)$ and
$(n_1, n_2) = (0.8636, 1.6840)$. The BHNS accretion-disk mass is the difference
\begin{equation}
\Mdisk = \max\!\left(0,\ M_{\rm rem} - M_{\rm dyn}\right).
\label{eq:disk}
\end{equation}
For a spin misaligned from the orbital angular momentum by a tilt angle
$\theta_{\rm tilt}$, only the aligned projection enters the fit through the
effective spin $\chi_{\rm eff} = \max(0,\ \aBH\cos\theta_{\rm tilt})$.

For completeness, the BNS post-merger disk admits a compactness-based fit; we
adopt a linearized form of the \citet{kruger2020} BNS disk relation (their
Eq.~4), $\Mdisk^{\rm BNS} = M_b^{\rm tot}\,\max(5\times10^{-4},\ a\,C_1 + c)$ with
$(a, c) = (-8.1580, 1.2695)$ and $C_1$ the compactness of the lighter component.
This simplifies the published Eq.~4,
$M_1\,\max\{5\times10^{-4},\,(a\,C_1 + c)^{d}\}$ with
$(a, c, d) = (-8.1324, 1.4820, 1.7784)$, by dropping the exponent and refitting
the linear coefficients (the published fit zeros the disk at $C_1 \approx 0.182$,
ours at $C_1 \approx 0.156$). The BNS classification below uses total-mass
and mass-ratio thresholds rather than this disk mass, because the BNS fit is
knife-edge sensitive to the NS radius near $C_1 \approx 0.156$. Neutron-star radii
in the fiducial disk-mass chain use the heuristic of Eq.~\ref{eq:nsradius} at
$R_{1.4} \approx 12$ km; the APR4, SFHo, LS220, and DD2 tabulations of
\citet{bauswein2013,bauswein2020} and \citet{read2009} supply the EOS-dependent
critical masses (and EOS-anchored radii where an EOS is named explicitly).

\subsection{GRB classification} \label{sec:classification}

We adopt the \citet{gottlieb2024} four-class hybrid scheme for BNS. With
$\Mtov = 2.2\,\Msun$ \citep{raaijmakers2021}, the prompt-collapse and
hyper-massive-neutron-star (HMNS) boundaries are
\begin{align}
\Mthresh &= k\,\Mtov = 1.27 \times 2.2\,\Msun = 2.794\,\Msun, \nonumber \\
M_{\rm HMNS} &= 1.2\,\Mtov = 2.64\,\Msun,
\label{eq:thresholds}
\end{align}
at the \citet{gottlieb2023} fiducial $k = 1.27$, with mass ratio
$q \equiv m_1/m_2 \geq 1$ and $q_{\rm thresh} = 1.2$. The four classes are
then defined on the total mass $\Mtot = m_1 + m_2$ and $q$; the full decision
tree is summarized in Figure~\ref{fig:scheme}.
The prompt-collapse threshold concept is from \citet{bauswein2013}, who give an
EOS-dependent $k \in [1.30, 1.70]$; the value $k = 1.27$ is a \citet{gottlieb2023}
fiducial chosen so that $\Mthresh \approx 2.8\,\Msun$. The long-lived HMNS split at
$M_{\rm HMNS} = 1.2\,\Mtov$ is a code heuristic motivated by the maximum mass that
uniform rotation can support \citep{margalit2017}, not a value tabulated by
\citet{gottlieb2024}. The equation of state enters only through $\Mtov$ and this
fiducial $k$; we adopt a single EOS-motivated $(\Mtov, k)$ pair throughout rather
than sweeping the EOS. Table~\ref{tab:params} collects the key physical
parameters of the classification, the rate corrections of
Section~\ref{sec:corrections}, and the observed-sample engine comparison of
Section~\ref{sec:obsmethod}, separating each fiducial value from the range
swept or marginalized.

\begin{deluxetable*}{llccl}
\tabletypesize{\footnotesize}
\tablecaption{Key physical parameters. \label{tab:params}}
\tablehead{
\colhead{Parameter} & \colhead{Symbol} & \colhead{Fiducial} & \colhead{Range} & \colhead{Source}
}
\startdata
\cutinhead{Classification (Sections~\ref{sec:remnant} and \ref{sec:classification})}
Maximum non-rotating NS mass     & $\Mtov$ & $2.2\,\Msun$ & \nodata & \citet{raaijmakers2021} \\
NS radius anchor                 & $R_{1.4}$ & $12$ km & \nodata & \citet{raaijmakers2021}; Eq.~\ref{eq:nsradius} \\
Prompt-collapse ratio            & $k$ & $1.27$ & \nodata & \citet{gottlieb2023} fiducial \\
Prompt-collapse threshold        & $\Mthresh = k\,\Mtov$ & $2.794\,\Msun$ & \nodata & derived \\
HMNS split                       & $M_{\rm HMNS} = 1.2\,\Mtov$ & $2.64\,\Msun$ & \nodata & \citet{margalit2017} heuristic \\
BNS mass-ratio threshold         & $q_{\rm thresh}$ & $1.2$ & \nodata & \citet{gottlieb2023} \\
BHNS jet-launch disk threshold   & $M_{\rm disk,short}$ & $0.01\,\Msun$ & \nodata & \citet{gottlieb2023} \\
BHNS bright-lbGRB disk threshold & $M_{\rm disk,long}$ & $0.1\,\Msun$ & \nodata & \citet{gottlieb2023} \\
\cutinhead{Rate corrections (Section~\ref{sec:corrections})}
sbGRB jet half-opening angle     & $\theta_j$ & $13^\circ$ & $10^\circ$ to $16^\circ$ & fiducial band; \citet{fong2015,beniamini_nakar2019} \\
lbGRB jet half-opening angle     & $\theta_j$ & $6.5^\circ$ & $5^\circ$ to $8^\circ$ & \citet{gottlieb2023} \\
BHNS black-hole spin             & $\aBH$ & $0.5$ & $\{0.0, 0.3, 0.5, 0.7, 0.9\}$ & this work; \citet{foucart2018} calibration \\
BHNS misalignment factor         & \nodata & $0.5$ & \nodata & \citet{fragos2010,kawaguchi2015} heuristic \\
\cutinhead{Observed-sample engine comparison (Section~\ref{sec:obsmethod})}
Jet-structure exponent           & $\alpha$ & \nodata & $1.5$ to $2$ & \citet{gottlieb2024} \\
Beaming-to-efficiency ratio      & $f_b/\epsilon_\gamma$ & \nodata & $0.02$ to $0.71$ & \citet{berger2014,beniamini2016} \\
HMNS disk-wind fraction          & $f_{\rm wind}$ & $0.3$ & $0.2$ to $0.4$ & \citet{radice2018} \\
HMNS-engine disk-mass prior      & $\Mdisk$ & \nodata & $0.01$ to $0.10\,\Msun$ & \citet{gottlieb2024} \\
\enddata
\tablecomments{Parameters with no Range entry are held fixed; the equation of
state enters only through the $(\Mtov, k)$ pair
(Section~\ref{sec:classification}). The NS radius anchor sets the mass-radius
relation of Eq.~\ref{eq:nsradius}, a modeling choice consistent with
\citet{raaijmakers2021} rather than an EOS derivation. The two BHNS disk
thresholds keep the \citet{gottlieb2023} short/long symbols but, under the
\citet{gottlieb2024} hybrid relabel of Figure~\ref{fig:scheme}, bound
the No GRB / faint lbGRB and faint lbGRB / lbGRB plus red kilonova (disk)
classes. The fiducial $\aBH = 0.5$ is a choice of this work, not a
\citet{foucart2018} recommendation; their remnant fit is calibrated for
$\chi_{\rm BH} \in [-0.5, 0.9]$, so spins are clipped at $0.9$ and the sweep
grid is restricted to prograde values. The sbGRB jet-angle band is the
fiducial compromise of Section~\ref{sec:corrections} (\citealt{fong2015}
measure a median $16^\circ \pm 10^\circ$; \citealt{beniamini_nakar2019}
prefer cores near $10^\circ$). Engine-comparison parameters with no fiducial
are marginalized over their Range; the disk-mass and wind-fraction priors are
log-uniform (Section~\ref{sec:obsmethod}).}
\end{deluxetable*}

For BHNS we use three disk-mass classes with thresholds $0.01$ and $0.1\,\Msun$
\citep{gottlieb2023}, shown in the right branch of Figure~\ref{fig:scheme}. The
class names follow the \citet{gottlieb2024} hybrid so they align with the BNS
labels. A unified mass-plane classifier assigns integer labels on the
$(m_1, m_2)$ plane for the diagnostic figures, truncating the BNS region at the
model-specific $M_{\rm NS,max}$.

\begin{figure*}[tbp]
\epsscale{1.05}
\plotone{fig15_classification_scheme.pdf}
\caption{GRB engine classification scheme \citep{gottlieb2023,gottlieb2024}.
BNS mergers (left branch) are classified on the total mass $\Mtot = m_1 + m_2$
and mass ratio $q \equiv m_1/m_2 \geq 1$, with the boundaries of
Eq.~\ref{eq:thresholds} ($M_{\rm HMNS} = 2.64\,\Msun$,
$\Mthresh = 2.794\,\Msun$) marked as ticks and the prompt-collapse bin forking
on $q_{\rm thresh} = 1.2$. BHNS mergers (right branch) are classified on the
accretion-disk mass $\Mdisk = M_{\rm rem} - M_{\rm dyn}$ (Eq.~\ref{eq:disk})
with thresholds at $0.01$ and $0.1\,\Msun$. Box colors follow the class
palette used throughout; the italic label below each box is the central
engine. The two confirmed BNS mergers are placed by the classifier at their
published median masses: GW170817 in the lbGRB plus red kilonova (HMNS) class
and GW190425 in the faint lbGRB class (Section~\ref{sec:massplane}). The
diagram is a schematic of the classification criteria; no \compas\ data
enter. \label{fig:scheme}}
\end{figure*}

\subsection{Formation channels} \label{sec:channels}

We decompose each population onto the five \citet{broekgaarden2021} formation
channels: I (stable mass transfer followed by a common envelope), II (stable mass
transfer only), III (single-core common envelope), IV (double-core common
envelope), and V (all others). Channel V is an exhaustive catch-all, so the
channel sum reproduces the full population to machine precision.

\subsection{Cosmic-rate integration} \label{sec:cosmicrate}

Intrinsic merger rates are computed with a chunked accumulator around the
\compas\ \fci\ routines. The intrinsic merger-rate density at merger redshift
$z_m$ is the delay-time convolution of the metallicity-specific star-formation
history with the per-system formation efficiency,
\begin{equation}
\mathcal{R}(z_m) = \sum_i \frac{w_i}{p_{\rm draw}\,\langle M_{\rm evolved}\rangle}\,
\mathcal{S}\!\left(Z_i, z_f\right),
\label{eq:rate}
\end{equation}
following \citet[their Eq.~1]{neijssel2019}, where the sum runs over \compas\
systems $i$ with STROOPWAFEL weight $w_i$ and birth metallicity $Z_i$, the
formation redshift $z_f$ is fixed by the system delay time $t_{d,i}$ through
$t(z_f) = t(z_m) - t_{d,i}$, and $p_{\rm draw}$ is the metallicity sampling
prior. The normalization $\langle M_{\rm evolved}\rangle$ is the mean stellar
mass evolved per \compas\ binary (the star-forming mass per simulated system;
\citealt{neijssel2019}, their Sec.~2.1). The
metallicity-specific star-formation rate density factorizes as
$\mathcal{S}(Z, z) = \mathrm{SFR}(z)\,dP/d\ln Z$. The cosmic star-formation rate
uses the \citet{madau2014} functional form
\begin{equation}
\mathrm{SFR}(z) = a\,\frac{(1+z)^b}{1 + [(1+z)/c]^d},
\label{eq:sfr}
\end{equation}
with the IllustrisTNG TNG100-1 best fit $(a, b, c, d) = (0.0172, 1.4425, 4.5299,
6.2261)$ \citep{levina2026}. The metallicity distribution is the
\citet{levina2026} Azzalini skew-log-normal, whose location and scale evolve as
\begin{equation}
\langle Z\rangle(z) = \mu_0\,10^{\mu_z z}, \qquad \omega(z) = \omega_0\,10^{\omega_z z},
\label{eq:mssfr}
\end{equation}
with skewness $\alpha$ and best fit $(\mu_0, \mu_z, \omega_0, \omega_z, \alpha) =
(0.0247, -0.0521, 1.1509, 0.0477, -1.8801)$; $dP/d\ln Z$ is normalized per
natural-log metallicity interval, which we verify numerically. We adopt the
Planck 2015 cosmology of the underlying \fci\
($H_0 = 67.74\,{\rm km\,s^{-1}\,Mpc^{-1}}$, $\Omega_m = 0.3089$,
$\Omega_\Lambda = 0.6911$; \citealt{planck2016}) throughout, to avoid a
high-redshift inconsistency in the rate integral. The per-population
$\langle M_{\rm evolved}\rangle$ is back-derived from the \compas\ local-rate
anchor encoded in the published $w_{000}$ weights using the \citet{neijssel2019}
MSSFR; this anchor is a simulation property and is MSSFR-independent, so the
Levina TNG100-1 MSSFR is used for all reported rates. Rates are intrinsic,
comoving, source-frame in $\Gpcyr$ unless stated otherwise. Because the disk-mass
classification assumes 100 percent jet launching above threshold, all GRB rates
are upper bounds \citep{gottlieb2023}.

\subsection{Rate corrections} \label{sec:corrections}

Three corrections enter specific results. (i) Beaming: to compare an intrinsic
rate with an observed sGRB rate we multiply by the beaming fraction
\begin{equation}
f_{\rm beam} = 1 - \cos\theta_j, \qquad \mathcal{R}_{\rm obs} = f_{\rm beam}\,\mathcal{R},
\label{eq:beam}
\end{equation}
with per-class jet half-opening angles $\theta_j = 13^\circ$ (range 10 to 16) for
the sbGRB class and $\theta_j = 6.5^\circ$ (range 5 to 8) for the lbGRB classes.
The sbGRB band is a fiducial prior rather than a single measurement: \citet{fong2015}
find a median sGRB half-opening angle $\theta_j = 16^\circ \pm 10^\circ$, while the
\citet{beniamini_nakar2019} structured-jet analysis of GW170817 prefers narrower
cores closer to $10^\circ$, and we adopt $10^\circ$ to $16^\circ$ as the
compromise band.
(ii) BH spin: BHNS rates are computed on the grid
$\aBH \in \{0.0, 0.3, 0.5, 0.7, 0.9\}$ with a fiducial $\aBH = 0.5$, and
marginalized over a spin prior $p(\aBH)$,
\begin{equation}
\mathcal{R}_{\rm marg} = \int \mathcal{R}(\aBH)\,p(\aBH)\,d\aBH
\approx \sum_j \mathcal{R}\!\left(\aBH^{(j)}\right) p\!\left(\aBH^{(j)}\right),
\label{eq:spinmarg}
\end{equation}
evaluated on the grid $\{\aBH^{(j)}\}$ for both a flat prior and the
\citet{fuller2019} low-spin prior.
(iii) Spin-orbit misalignment: BHNS population rates are scaled by an
order-unity heuristic factor, $\mathcal{R}_{\rm BHNS} \to 0.5\,\mathcal{R}_{\rm BHNS}$,
motivated by the misaligned-spin distributions of \citet{fragos2010} together
with the numerical-relativity result that the remnant disk mass collapses for
tilt angles beyond $\sim 50$ to $60^\circ$ \citep{kawaguchi2015}; it is applied
at the rate level only, not by dropping systems. We do not apply beaming or the misalignment factor when comparing against
the GWTC-5.0 intrinsic GW rate, which already folds in selection effects.

\subsection{Observed-sample engine comparison} \label{sec:obsmethod}

To confront the engine branches with data we use the eight kilonova-bearing GRBs
of \citet{rastinejad2024}, the only sample with a uniform blue/purple/red
ejecta decomposition $(M_B, M_P, M_R)$ (their Table~3). For each event we draw
$10^4$ Monte Carlo realizations of the component ejecta masses from split-normal
posteriors and of the burst duration $T_{50}$ (with log-uniform or log-normal
priors for the two events lacking a published $T_{50}$). The black-hole-engine
prediction for the ejecta mass at a given duration follows the
\citet[their Eqs.~10 and 11]{gottlieb2024} scaling
$M_{\rm ej} \propto f^{-1}\,E_{\gamma,\rm iso}\,T_{50}^{\,\alpha-1}$, marginalized
over the jet-structure exponent $\alpha \in [1.5, 2]$ and the ratio
$f_b/\epsilon_\gamma \in [0.02, 0.71]$ set by the \citet{berger2014} beaming
and \citet{beniamini2016} radiative-efficiency bounds. The HMNS-engine
prediction is the thermal disk-wind ejecta $M_{\rm ej} = f_{\rm wind}\,\Mdisk$,
with $f_{\rm wind} \in [0.2, 0.4]$ spanning the GRMHD wind fractions of
\citet{radice2018} and $\Mdisk \in [0.01, 0.10]\,\Msun$
\citep[their Sec.~3.1]{gottlieb2024}. The diagnostic is the residual
$\log_{10}(M_{\rm obs,wind}/M_{\rm pred,BH})$, which is consistent with zero for a
black-hole-disk engine and strongly positive for an HMNS-wind engine. This
comparison is an external consistency check on the engine classes and does not
feed back into the population classification.

\subsection{Model assumptions and caveats} \label{sec:assumptions}

For transparency we collect here the assumptions that most affect the results,
each with its source.
\begin{enumerate}
\item \textit{Cosmology.} Planck 2015 \citep{planck2016}, matching the underlying
\fci, is used throughout; mixing in a different cosmology shifts the high-redshift
rate integral.
\item \textit{Star-formation history.} The MSSFR is the \citet{levina2026}
IllustrisTNG TNG100-1 fit. As that work shows, analytical $S(Z,z)$ prescriptions
can overestimate double-compact-object rates at high redshift by up to several
orders of magnitude relative to the simulation-based history; our absolute rates
inherit this systematic, which is why we anchor to the local rate and emphasize
relative (per-class, per-channel) results.
\item \textit{Importance weighting.} The \compas\ samples are STROOPWAFEL
importance-sampled; every reduction is weighted, and unweighted statistics would
be biased.
\item \textit{Neutron-star masses.} The Fryer mass gap is closed by remapping onto
the \citet{alsing2018} double Gaussian \citep{mandel2020}; this shifts class
fractions relative to the raw catalog and is applied before any classification.
\item \textit{Supernova engine.} All core-suite models use the \citet{fryer2012}
delayed mechanism; the rapid variation (Model I) enters only the grid scan of
Appendix~\ref{app:grid}.
\item \textit{Black-hole spin.} The \citet{foucart2018} remnant fit is calibrated
for $\chi_{\rm BH} \in [-0.5, 0.9]$; spins are clipped at $0.9$ to avoid
extrapolation.
\item \textit{Prompt-collapse threshold.} We adopt $\Mtov = 2.2\,\Msun$
\citep{raaijmakers2021} and the \citet{gottlieb2023} fiducial $k = 1.27$, below the
\citet{bauswein2013} EOS band $[1.30, 1.70]$; the equation of state enters only
through this pair.
\item \textit{Jet launching.} All disk-mass-based GRB rates assume 100 percent jet
launching above threshold \citep{gottlieb2023} and are therefore upper bounds.
\item \textit{BHNS spin-orbit misalignment.} A population-averaged, order-unity
heuristic factor $0.5$ (motivated by \citealt{fragos2010,kawaguchi2015}) is
applied at the rate level, not by removing systems, and is omitted from the
GWTC-5.0 comparison.
\item \textit{Quoted local rates.} All $\Rzero$ values are read from the
unsmoothed $\Rz$; the light Gaussian kernel that suppresses
discrete-metallicity-grid wobble in the plotted curves reflects at $z = 0$ and
would otherwise inflate the local rate by $\sim 30$ percent.
\end{enumerate}

%% ======================================================================
\FloatBarrier
\section{Results} \label{sec:results}
%% ======================================================================

Unless noted, results are for the fiducial Model A. The remapped, weighted BNS
sample contains $N = 233{,}137$ merging systems and the BHNS sample
$N = 1{,}525{,}553$ (classified at the fiducial $\aBH = 0.5$; the sample itself
is spin-independent). The headline class fractions and local
rates are collected in Table~\ref{tab:classfrac}: roughly three-quarters of BNS
mergers fall in the lbGRB-producing classes, and the BHNS population is dominated
by NS-swallowing systems that launch no GRB. The results below build on
these fractions in order: where the classes sit in the mass plane and what mass,
delay-time, and metallicity distributions they imply
(Sections~\ref{sec:massplane} to \ref{sec:delaytime}); the cosmic merger rate
and its decomposition by class and by black-hole spin
(Sections~\ref{sec:classrates} and \ref{sec:spin}; the formation-channel
decomposition is in Appendix~\ref{app:channels}); the comparison with
the GWTC-5.0 local rate (Section~\ref{sec:lvk}); and the confrontation with the
observed kilonova-bearing sample, from class fractions to host-galaxy offsets
(Sections~\ref{sec:obsclass} to \ref{sec:offsets}).

\begin{deluxetable*}{lccc}
\tablecaption{Model A class fractions and local merger rates by GRB class.
\label{tab:classfrac}}
\tablehead{
\colhead{Class} & \colhead{Raw $N$} & \colhead{Fraction} & \colhead{$\Rzero$} \\
\colhead{} & \colhead{} & \colhead{(weighted)} & \colhead{($\Gpcyr$)}
}
\startdata
\cutinhead{Binary neutron star ($N = 233{,}137$, $n_{\rm eff} = 65{,}134$)}
sbGRB plus blue kilonova       & $52{,}660$ & 24.0\% & 21.8 \\
lbGRB plus red kilonova (HMNS) & $87{,}852$ & 26.8\% & 19.7 \\
lbGRB plus red kilonova (disk) & $2{,}614$  & 2.2\%  & 1.4 \\
Faint lbGRB                    & $90{,}011$ & 47.0\% & 12.1 \\
All BNS                        & $233{,}137$ & 100\% & 55.0 \\
\cutinhead{Black hole neutron star ($N = 1{,}525{,}553$, $n_{\rm eff} = 292{,}484$)}
lbGRB plus red kilonova (disk) & $7{,}573$    & 3.8\%  & 3.9 \\
Faint lbGRB                    & $125{,}479$  & 20.6\% & 19.8 \\
No GRB (NS swallowed)          & $1{,}392{,}501$ & 75.6\% & 60.8 \\
All BHNS                       & $1{,}525{,}553$ & 100\% & 84.5 \\
\enddata
\tablecomments{\compas\ Model A, STROOPWAFEL-weighted; fractions are weighted,
$n_{\rm eff}$ is the Kish effective sample size. Rates are intrinsic, comoving,
source-frame local rate densities $\Rzero$ under Planck 2015 cosmology
\citep{planck2016}, read from the unsmoothed $\Rz$ (Section~\ref{sec:assumptions})
and therefore consistent with the per-model bars of
Figure~\ref{fig:rate20}; BHNS rates are intrinsic at $\aBH = 0.5$ without
the misalignment factor. All disk-mass-based GRB rates are upper bounds
\citep{gottlieb2023}. Per-class rates are rounded to the quoted precision and
need not sum exactly to the all-population total.}
\end{deluxetable*}

\subsection{The mass plane} \label{sec:massplane}

\begin{figure*}[tbp]
\epsscale{1.15}
\plottwo{mass_plane_bns.pdf}{mass_plane_bhns.pdf}
\caption{Compact-binary mass plane for \compas\ Model A
\citep[Zenodo 5189849 and 5178777]{broekgaarden2021}. \textit{Left:} BNS, with the
\citet{gottlieb2024} four-class hybrid boundaries ($\Mtot = 1.2\,\Mtov$,
$\Mtot = \Mthresh = 1.27\,\Mtov$, $q = q_{\rm thresh} = 1.2$) and a four-fold
reflective KDE; purple stars mark the two confirmed BNS mergers GW170817 and
GW190425. \textit{Right:} BHNS, with
\citet{foucart2018} disk-mass contours at $0.01$ and $0.1\,\Msun$ at
$\aBH = 0.5$. Samples are STROOPWAFEL-weighted: BNS raw $N = 233{,}137$,
$n_{\rm eff} = 65{,}134$; BHNS raw $N = 1{,}525{,}553$, $n_{\rm eff} = 292{,}484$.
NS masses are remapped onto the \citet{alsing2018} double Gaussian
\citep{mandel2020}; $\Mtov = 2.2\,\Msun$ \citep{raaijmakers2021}.
\label{fig:massplane}}
\end{figure*}

Figure~\ref{fig:massplane} shows the $(m_1, m_2)$ mass plane for Model A. The BNS
panel reproduces and extends \citet[their Fig.~3]{gottlieb2024} using a
four-fold reflective kernel density estimate (symmetrized under
$m_1 \leftrightarrow m_2$ and mirrored at $\Mtov$), with the hybrid classification
boundaries ($\Mtot = \Mthresh$, $\Mtot = 1.2\,\Mtov$, $q = 1.2$) overlaid and the
two confirmed BNS mergers GW170817 and GW190425 marked. GW170817 itself
($\Mtot \approx 2.73\,\Msun$; \citealt{abbott2019gw170817}) falls
$0.09\,\Msun$ above the $M_{\rm HMNS}$ boundary and is therefore assigned to
the lbGRB plus red kilonova (HMNS) class rather than to the sbGRB class one
might expect for GRB\,170817A. The assignment is boundary-sensitive: the
$1.2\,\Mtov$ split is a heuristic (Section~\ref{sec:classification}), and
$\Mtov \geq 2.28\,\Msun$, well within current EOS uncertainty, moves GW170817
across it. GW190425 ($\Mtot \approx 3.4\,\Msun$;
\citealt{abbott2020gw190425}) sits above $\Mthresh$ at near-unity mass ratio
and is assigned to the faint lbGRB (prompt-collapse) class, consistent with
its lack of a detected counterpart. Per-event assignments near the boundaries
inherit this EOS sensitivity even though the population fractions are
comparatively stable (Appendix~\ref{app:robustness}). The BHNS panel shows the
disk-mass classes at $\aBH = 0.5$ with \citet{foucart2018} disk-mass contours at
the $0.01$ and $0.1\,\Msun$ thresholds.

\subsection{Component-mass distributions by class} \label{sec:massdist}

\begin{figure*}[tbp]
\epsscale{1.15}
\plotone{mass_distributions_by_class.pdf}
\caption{Weighted, abundance-normalized distributions of $m_1$, $m_2$, $\Mtot$,
and $q$ by GRB class (top row BNS, bottom row BHNS) for \compas\ Model A. Areas
equal class fractions; the two BNS lbGRB plus red kilonova classes share color and
differ by line style (HMNS solid, disk dashed). STROOPWAFEL-weighted; samples as
in Figure~\ref{fig:massplane}. The weighted-Poisson per-bin uncertainty is below
the line width at these effective sample sizes, so the panels convey shape; the
absolute normalization enters with rate units in Section~\ref{sec:classrates}.
\label{fig:massdist}}
\end{figure*}

Figure~\ref{fig:massdist} shows weighted, abundance-normalized distributions of
$m_1$, $m_2$, $\Mtot$, and $q$ split by GRB class for both populations. The
curves within each row jointly integrate to unity, so the area under each class
curve equals its population fraction. The two BNS lbGRB plus red kilonova classes
share the project red color and are distinguished by line style. The sbGRB plus
blue kilonova class occupies the low-$\Mtot$ tail, consistent with its
$\Mtot < 1.2\,\Mtov$ definition, while the faint lbGRB class dominates at high
$\Mtot$ and low $q$. Because the classifier operates on the remapped masses
(Section~\ref{sec:remap}), these abundance fractions
(Table~\ref{tab:classfrac}) inherit the \citet{alsing2018} double-Gaussian
target distribution \citep{mandel2020}; a different mass remap would move the
boundaries between the low-$\Mtot$ and high-$\Mtot$ classes.

\subsection{Delay-time and metallicity distributions} \label{sec:delaytime}

\begin{figure*}[tbp]
\epsscale{1.15}
\plotone{delay_time_distributions.pdf}
\par\medskip
\plotone{metallicity_dependence.pdf}
\caption{Per-class formation properties for \compas\ Model A,
STROOPWAFEL-weighted (BNS left, BHNS right). \textit{Top:} weighted delay-time
distributions by GRB class, per-class normalized over
$t_{\rm delay} \in [3, 1.4\times10^4]$ Myr. \textit{Bottom:} weighted formation
efficiency $\eta(Z)$ per GRB class versus $Z/\Zsun$, with $\Zsun = 0.0142$
\citep{asplund2009}; bins require raw $N \geq 50$, Kish $n_{\rm eff} \geq 10$,
and $Z \leq \Zsun$. The weighted-Poisson per-bin uncertainty is below the line
width at these effective sample sizes (the coarse 30-bin logarithmic spacing of
the delay-time panels further suppresses single-system spikes).
\label{fig:delaytime}}
\end{figure*}

The top row of Figure~\ref{fig:delaytime} shows weighted delay-time
distributions $t_{\rm delay} = t_{\rm form} + t_c$ by class on logarithmic axes,
normalized per class over the visible range so that the panels show distribution
shape rather than abundance (abundance enters the rate in
Section~\ref{sec:classrates}). All classes follow the expected approximately
$t^{-1}$ behavior at late times; the faint lbGRB and disk classes extend to the
shortest delays, while the sbGRB plus blue kilonova class switches on latest,
near 30 Myr.

The bottom row shows the intensive formation efficiency
$\eta(Z_i) = (\sum w)_i / M_{\rm SF}(Z_i)$ per class, the number of merging systems
formed per unit star-forming mass at each metallicity \citep[their Fig.~5]{broekgaarden2021}.
The per-metallicity star-forming mass is the population total split by the \compas\
uniform-in-$\ln Z$ sampling prior, $M_{\rm SF}(Z_i) = \langle M_{\rm evolved}\rangle\,f_i$,
where $f_i$ is the prior mass fraction in bin $i$; dividing by this per-$Z$ mass rather than
the scalar total makes $\eta(Z)$ intensive (independent of how the grid is binned). The
53-point \compas\ metallicity grid is pooled into equal-count logarithmic bins for display
only. Bins are retained when the raw count is at least 50, the Kish effective size at least
10, and $Z \leq \Zsun = 0.0142$ \citep{asplund2009}. The disk-mass and faint
classes rise toward low metallicity, reflecting the metallicity dependence of
compact-object formation in \citet{broekgaarden2021}.

\subsection{Merger rate by GRB class} \label{sec:classrates}

\begin{figure*}[tbp]
\epsscale{1.15}
\plotone{rate_bns_by_class.pdf}
\par\medskip
\plotone{rate_bns_tng_sweep.pdf}
\caption{Intrinsic BNS merger rate $\Rz$ for \compas\ Model A. \textit{Top:}
$\Rz$ per \citet{gottlieb2024} GRB class together with the all-BNS total (HMNS
solid, disk dashed, shared color). \textit{Bottom:} the all-BNS $\Rz$ under the
three \citet{levina2026} IllustrisTNG MSSFR resolutions (TNG50-1, TNG100-1,
TNG300-1), with the resolution envelope shaded; this envelope is the dominant
systematic on the per-class curves and is common to all classes, while the
weighted-Poisson sampling error is sub-percent. Cosmology: Planck 2015
\citep{planck2016}; rates are intrinsic, comoving, source-frame, and the
GRB-class rates are upper bounds \citep{gottlieb2023}. \label{fig:rateclass}}
\end{figure*}

The remaining results convert these distributions into cosmic rates.
Figure~\ref{fig:rateclass} (top) shows the intrinsic BNS $\Rz$ per GRB class
together with the all-BNS total to $z = 10$. The four-class partition sums to
the total at the level of the numerical tolerance. The local-rate ordering
inverts relative to the abundance ordering of Section~\ref{sec:massdist}: the
sbGRB plus blue kilonova and HMNS classes carry the bulk of $\Rzero$ ($21.8$
and $19.7\,\Gpcyr$; Table~\ref{tab:classfrac}), while the faint lbGRB class, 47
percent of the population by weight, contributes $12.1\,\Gpcyr$ (22 percent of
the local rate). The inversion follows from the per-class delay-time and
metallicity distributions: the faint class forms preferentially at low
metallicity and short delay (Section~\ref{sec:delaytime}), so its mergers
concentrate at higher redshift and are relatively depleted at $z = 0$.

The dominant systematic on these curves is the star-formation history. The
bottom panel of Figure~\ref{fig:rateclass} repeats the all-BNS $\Rz$ for the
three IllustrisTNG resolutions exposed by \citet{levina2026}, TNG50-1,
TNG100-1, and TNG300-1, with TNG100-1 as the fiducial. The published Levina BBH
local rates under these resolutions, $73.7$, $45.5$, and $27.8\,\Gpcyr$
respectively, set the amplitude ordering
$R_{\rm TNG50} > R_{\rm TNG100} > R_{\rm TNG300}$ that our BNS curves inherit,
driven by the star-formation normalization and the skew-PDF amplitude rather
than the compact-object type. The shaded band is the resolution envelope; it
reappears as the whisker on the all-BNS bar of Figure~\ref{fig:lvkbars}.

The decomposition of these per-class rates onto the \citet{broekgaarden2021}
formation channels of Section~\ref{sec:channels} is collected in
Appendix~\ref{app:channels}.

\subsection{BHNS BH-spin sensitivity} \label{sec:spin}

\begin{figure*}[tbp]
\epsscale{1.15}
\plotone{rate_bhns_spin_sensitivity.pdf}
\caption{BHNS intrinsic merger rate $\Rz$ versus BH spin for \compas\ Model A.
Shaded envelopes span $\aBH \in \{0.0, 0.3, 0.5, 0.7, 0.9\}$, the
\citet{foucart2018} calibration band; the central line is the fiducial
$\aBH = 0.5$. Per-class curves include the misalignment factor 0.5; the
all-BHNS intrinsic reference is shown without it. Cosmology: Planck 2015
\citep{planck2016}; GRB-class rates are upper bounds \citep{gottlieb2023}. The
shaded envelopes are themselves the BH-spin systematic, the dominant
within-model uncertainty on the BHNS class rates. \label{fig:spinsweep}}
\end{figure*}

Figure~\ref{fig:spinsweep} sweeps $\aBH$ over the \citet{foucart2018}
calibration band and reclassifies the BHNS population at each spin. Higher spin
shrinks the innermost stable circular orbit, increases the disk mass, and moves
systems from the no-GRB class into the GRB-launching classes, so both the lbGRB
and faint lbGRB rates rise monotonically with spin. The spread across the spin
grid is the envelope carried forward as the per-class systematic in
Figure~\ref{fig:lvkbars}.

\subsection{Local rate versus LVK GWTC-5.0} \label{sec:lvk}

We compare the Model A local rate densities $\Rzero$ with the
GWTC-5.0 90 percent credible regions \citep{abac2026}: the joint BNS region
$[5.1, 154.7]\,\Gpcyr$ and the joint NSBH region $[6.7, 32.8]\,\Gpcyr$ (both at
$z = 0$, the union of the PixelPop and FullPop population models). Both the model
and the LVK measurement are intrinsic source-frame rates, so no beaming or
misalignment correction is applied here. The GWTC-5.0 per-model BNS estimates are
$23.4^{+54.7}_{-18.2}$ (PixelPop) and $59.3^{+95.4}_{-43.9}$ (FullPop), and the NSBH
estimates are $15.9^{+16.9}_{-8.9}$ and
$14.2^{+12.0}_{-7.5}\,\Gpcyr$. GWTC-5.0 extends the sample through the second part
of the fourth observing run but adds no new BNS or NSBH detections, so the
neutron-star regions tightened relative to GWTC-4 (which gave BNS
$[7.6, 250]$ and NSBH $[9.1, 84]\,\Gpcyr$; \citealt{abac2025}). The BNS rate,
$55\,\Gpcyr$, falls comfortably within the GWTC-5.0 BNS region and within the
FullPop per-model interval. The BHNS rate, $84.5\,\Gpcyr$, sits a factor
$\sim 2.6$ above the NSBH 90 percent upper bound and well above both per-model
medians; Section~\ref{sec:discussion} discusses what this implies. Across the
five population-synthesis variations every BHNS rate exceeds the tightened
NSBH region, by factors of $\sim 1.2$ to $2.9$ (per-model rates in
Appendix~\ref{app:robustness}; Figure~\ref{fig:rate20}). Independent
multi-messenger analyses point the same way: the GW-inferred NS merger rates
are low enough to strain the observed sGRB rate at canonical jet angles,
implying either genuinely low rates \citep{fishbach2026,jin2026} or wider
short-GRB jets \citep{kunnumkai2026}; Section~\ref{sec:priorwork} returns to
this tension.

\begin{figure*}[tbp]
\epsscale{1.15}
\plottwo{rate_bns_vs_lvk_gwtc4.pdf}{rate_bhns_vs_lvk_gwtc4.pdf}
\caption{Local intrinsic merger rate $\Rzero$ by GRB class against the GWTC-5.0
90 percent credible regions \citep[shaded]{abac2026}. \textit{Left:} BNS, with
the hatched all-BNS bar (horizontal whisker: TNG50-1 to TNG300-1 MSSFR-resolution
systematic) and the four \citet{gottlieb2024} class subsets that sum to it.
\textit{Right:} BHNS at $\aBH = 0.5$, with the hatched all-BHNS bar and the two
GRB-class subsets; the class bars carry the BH-spin envelope over
$\aBH \in [0,0.9]$. Both axes are intrinsic source-frame rates, so no beaming or
misalignment correction is applied; all GRB-class rates are upper bounds
\citep{gottlieb2023}. \label{fig:lvkbars}}
\end{figure*}

Figure~\ref{fig:lvkbars} resolves the same comparison by GRB class.
The all-population bar is shown against the GWTC-5.0 joint region, and the
four BNS (three BHNS) class subsets that partition it are stacked below. The
all-BNS bar carries a horizontal whisker spanning the MSSFR-resolution
systematic (the TNG50-1 to TNG300-1 spread of Section~\ref{sec:classrates}),
which is the dominant model uncertainty on the local rate. For BHNS the two
GRB-class bars carry the black-hole-spin envelope over $\aBH \in [0,0.9]$, which
is the dominant per-class systematic (Section~\ref{sec:spin}); the all-BHNS bar
is spin-independent. Even across the spin envelope the GRB-producing BHNS classes
remain a small fraction of the all-BHNS total, so the BHNS rate that overshoots
the NSBH region is carried by the NS-swallowing no-GRB systems rather than by the
kilonova-bearing population.

\subsection{Observed class fractions} \label{sec:obsclass}

\begin{figure}[tbp]
\epsscale{1.15}
\plotone{obs_class_fractions.pdf}
\caption{Predicted versus observed \citet{gottlieb2024} class fractions. Filled
bars: abundance-weighted \compas\ Model A BNS fractions (Section~\ref{sec:massdist};
Alsing-remapped masses). Hatched bars: the eight kilonova-bearing GRBs of
\citet{rastinejad2024} (raw $N = 8$), mapped onto the four classes
through the blue/purple/red ejecta decomposition with the heuristic
$f_{\rm red}$ and $M_{\rm ej}$ thresholds described in the text. Observed error
bars are per-class multinomial uncertainties; the model bars carry no
sampling-error whisker (the weighted-Poisson error is sub-percent, and a
bootstrap treatment is deferred). The mapping is a consistency
check, not a fit. \label{fig:obsclass}}
\end{figure}

Figure~\ref{fig:obsclass} confronts the predicted Model A class fractions with the
eight kilonova-bearing GRBs of \citet{rastinejad2024} introduced in
Section~\ref{sec:obsmethod}. We map each
observed event onto the \citet{gottlieb2024} classes through its ejecta
decomposition. Writing the red fraction as
$f_{\rm red} = M_R / (M_B + M_P + M_R)$, a source is faint lbGRB when the total
ejecta mass is $M_{\rm ej} < 0.01\,\Msun$ (this cut takes precedence);
otherwise it is sbGRB plus blue kilonova when $f_{\rm red} < 0.30$, lbGRB plus
red kilonova (disk) when $f_{\rm red} \geq 0.70$, and lbGRB plus red
kilonova (HMNS) in between. These three thresholds partition the color plane; they
are code heuristics, not quantities tabulated by \citet{gottlieb2024} or
\citet{rastinejad2024}, and the comparison is a consistency check, not a fit.
The model bars are the abundance-weighted Model A fractions of
Table~\ref{tab:classfrac}, not the cosmic-rate-weighted fractions of
Section~\ref{sec:classrates}.

With only eight events the observed fractions carry large multinomial uncertainty,
so the comparison is qualitative. We attach to each observed fraction $f_k$ the
normal-approximation multinomial error $\sigma_k = \sqrt{f_k(1 - f_k)/N}$ with
$N = 8$. The observed sample splits into two sbGRB plus blue,
four mixed (HMNS), two red (disk), and zero faint. The model reproduces the
sbGRB-minority, long-duration-majority ordering but predicts a large
faint-lbGRB fraction with no observed counterpart. The absence is expected:
faint, ejecta-poor events are the least likely to yield a detected kilonova,
so a kilonova-selected sample excludes them by construction. A larger
uniformly modeled sample is required before the faint fraction can be tested.

\subsection{Durations and ejecta of the observed sample} \label{sec:enginediscriminator}

\begin{figure*}[tbp]
\epsscale{1.15}
\plottwo{comparison_fig1.pdf}{comparison_fig2.pdf}
\caption{Engine discriminator for the \citet{rastinejad2024} kilonova-bearing
GRBs. \textit{Left:} duration $T_{50}$ versus total kilonova ejecta mass. The
gray band is the black-hole-engine prediction over $\alpha \in [1.5, 2]$
anchored at the in-sample lbGRB mean $E_{\gamma,\rm iso}$, the dashed envelope
adds the $f_b/\epsilon_\gamma$ range, and the green strip is the HMNS disk-wind
$M_{\rm ej} = 0.3\,\Mdisk$ prior (\citealt{gottlieb2024}, their Sec.~3.1); error
bars are 68 percent on $M_{\rm ej,tot}$ and $T_{50}$. \textit{Right:} residual
$\log_{10}(M_{\rm obs,wind}/M_{\rm pred,BH})$ versus $T_{50}$, with the
BH-engine null at zero (gray) and the pooled HMNS-engine expectation bands at 68
percent. Per-event residuals span the Monte Carlo 16th to 84th percentiles,
marginalized over the disk-mass and wind-fraction priors. Asterisked events (open markers in the
right panel) use a synthetic
$T_{50}$ prior, lacking a published duration. \label{fig:enginef1}}
\end{figure*}

The class fractions of Section~\ref{sec:obsclass} use only the kilonova color;
the duration and ejecta mass of each burst carry independent engine information.
Figure~\ref{fig:enginef1} places the \citet{rastinejad2024} events in the
plane of burst duration and ejecta mass against the two engine predictions of
Section~\ref{sec:obsmethod} (left) and collapses the comparison onto the
residual $\log_{10}(M_{\rm obs,wind}/M_{\rm pred,BH})$ versus $T_{50}$ (right).
The result inverts the naive duration-to-engine mapping. The
long-duration bursts with kilonovae (060614, 211211A, 230307A) land on the
black-hole accretion-disk prediction, with a pooled offset of $+0.10$ dex
($+0.3\sigma$) from the BH-engine null, while exceeding the HMNS disk-wind
expectation by $+1.0$ dex ($+4.2\sigma$). The short-duration bursts (050709,
130603B, 160821B, 200522A) show the opposite preference: they sit $+3.4$ dex
($+12\sigma$) above the BH-disk prediction and come within $+1.0$ dex
($+4.7\sigma$) of the HMNS band, so their ejecta demand the HMNS-wind
enrichment even though the pooled offset does not close entirely. GW170817
follows the short-burst pattern when treated as on-axis, sitting at the upper
edge of the HMNS band; its face-value off-axis energetics place it above every
prediction. This is precisely the \citet{gottlieb2024} picture in which the central
engine, not the observed duration class, sets the kilonova: the bursts whose
ejecta are consistent with a black-hole disk are the long-duration ones, and the
short bursts require the magnetar or HMNS-wind enrichment. With seven events the
trend is suggestive rather than decisive, and the residual folds in the wide
log-uniform priors on the disk mass and wind fraction
(Section~\ref{sec:obsmethod}); it is nonetheless the first quantitative
confrontation of the Gottlieb two-engine taxonomy with a uniformly modeled
kilonova sample.

\subsection{Host-galaxy offsets} \label{sec:offsets}

\begin{figure*}[tbp]
\epsscale{1.15}
\plotone{projected_offsets.pdf}
\caption{Weighted projected galactocentric offset CDFs by GRB class for \compas\
Model A (BNS left, BHNS right), from \citet{hernquist1990} orbit integration with
\compas\ systemic velocities and a mixed star-forming/elliptical host
population. Observed comparators: the \citet{fong2013} sGRB sample (left) and the
lGRB-plus-kilonova offsets of \citet{dellavalle2006,rastinejad2022,levan2024}
(right). The
annotated $p$-value is a weighted two-sample KS statistic. STROOPWAFEL-weighted.
\label{fig:offsets}}
\end{figure*}

The two engine branches share the same compact-binary kinematics, so the
projected galactocentric offset distribution is a cross-check rather than a
discriminator, but it is the most direct contact with the localized sGRB and
lGRB-plus-kilonova samples. We integrate each system on a \citet{hernquist1990}
host potential with the post-supernova systemic velocity from \compas, drawing
birth radii from the Hernquist profile and assigning a mixed
star-forming/elliptical host population, and build the weighted offset
cumulative
distribution per GRB class (Figure~\ref{fig:offsets}). Weighted two-sample
Kolmogorov-Smirnov tests against the \citet{fong2013} sGRB offsets give no
significant difference for the BNS classes ($p = 0.17$ for the annotated lbGRB
plus red kilonova HMNS class), and the model BNS offsets reproduce the
observed sub-10-kpc median. The BHNS offsets extend slightly further; the
driver is the delay-time distribution rather than the kicks, since the
weighted-median BHNS systemic velocity ($82\,{\rm km\,s^{-1}}$) is below the
BNS value ($113\,{\rm km\,s^{-1}}$) while the BHNS delay times are an order of
magnitude longer (weighted medians $1.3$ Gyr versus $0.13$ Gyr). The
prediction is also insensitive to the natal-kick prescription itself
(Figure~\ref{fig:offsetkicks}, Appendix~\ref{app:robustness}).

%% ======================================================================
\FloatBarrier
\section{Discussion} \label{sec:discussion}
%% ======================================================================

The class fractions place most BNS mergers in the lbGRB-producing classes (HMNS
plus disk plus faint, about 76 percent by weight), with a substantial sbGRB plus
blue kilonova minority (24 percent). This is a direct consequence of the remapped
total-mass distribution sitting near the $1.2\,\Mtov$ and $\Mthresh$ boundaries:
the neutron-star mass remap of Section~\ref{sec:remap} raises component masses to
fill the Fryer mass gap, pushing systems across the sbGRB boundary and out of the
lowest-mass class. The BHNS population is dominated by NS-swallowing systems with
no GRB (76 percent), so the GRB-relevant BHNS rate is a small fraction of the
total merger rate even before beaming.

The comparison with GWTC-5.0 (Section~\ref{sec:lvk}) is the sharpest external
check. The BNS rate is consistent with the broad GWTC-5.0 BNS credible region;
the BHNS rate is not, for the fiducial model or for any other member of the
core suite. What the comparison probes is \compas\ BHNS production at the
calibrated local-rate normalization, so we read the overshoot as a generic
overprediction by the core-suite physics, sharpest for the highest maximum NS
mass (Model K) but not isolated to $M_{\rm NS,max}$. The full grid shows which
assumptions could relieve it: fixed mass-transfer efficiency ($\beta = 0.5$ or
$0.75$), $\alpha_{\rm CE} = 10$, or enhanced Wolf-Rayet winds pull the BHNS
rate inside the NSBH region while keeping the BNS rate inside its own,
whereas unstable case-BB mass transfer does so only by driving the BNS rate
below its band (Section~\ref{sec:systematics}; Figure~\ref{fig:rate20}).

\citet{broekgaarden2026} give the BNS side of this comparison a direct
physical reading. Compiling formation-pathway fractions across
population-synthesis codes, they find that merging BNS formation is
overwhelmingly mediated by common-envelope evolution, so the observed BNS
merger rate directly traces the efficiency and success of envelope ejection;
a deficit of BNS mergers relative to the BBH and BHNS populations would point
to inefficient ejection or a high rate of failed CE events. Our grid is
consistent with the premise: the only CE-free pathway, channel II (stable
mass transfer only), never reaches half a percent of the merging BNS
population in the five-model suite (Figure~\ref{fig:chanfrac}), and the BNS
local rate responds directly to the efficiency, moving from $29.3$ (Model F,
$\alpha_{\rm CE} = 0.5$) to $106.0\,\Gpcyr$ (Model G, $\alpha_{\rm CE} = 2.0$)
while the class fractions barely move (Appendix~\ref{app:grid}). In this
framework the absolute BNS rate is a common-envelope anchor, and the class
demographics are the robust observable layered on top of it.

The same reading sharpens what is at stake in the evolving GW constraint.
\citet{broekgaarden2026} note that recent population analyses suggest an
intrinsic BNS rate lower than the rates inferred from the first detections,
potentially in tension with the short-GRB, Galactic double-neutron-star, and
r-process constraints \citep{fishbach2026,abac2026}, although
\citet{kunnumkai2026} read the same comparison as favoring wider short-GRB
jets. Should the lower values hold, our fiducial $55\,\Gpcyr$ sits on the
high side, and within the \citet{broekgaarden2021} grid the revision would
push toward less efficient envelope ejection. No equivalent leverage exists
on the BHNS side, where \citet{broekgaarden2026} find the relative importance
of the CE and stable mass-transfer channels remains highly uncertain, so the
BHNS overprediction of Section~\ref{sec:lvk} constrains the overall \compas\
BHNS production rather than the envelope physics.

The robustness study (Appendix~\ref{app:robustness}) shows the qualitative class
ordering is stable against the common-envelope efficiency and the maximum NS mass,
while the detailed fractions shift at the tens-of-percent level. The BH-spin
sweep (Figure~\ref{fig:spinsweep}) is the largest single source of BHNS rate
uncertainty within a fixed model, since the disk mass and therefore the GRB
classification depend strongly on spin through the innermost stable circular
orbit.

\subsection{Comparison with previous work} \label{sec:priorwork}

Our intrinsic BHNS local rate, $84.5\,\Gpcyr$ at $\aBH = 0.5$, falls inside the
broad COMPAS BHNS range of $4$ to $830\,\Gpcyr$ that \citet{broekgaarden2021}
obtain across their 420 model permutations; its position above the GWTC-5.0
NSBH upper bound (Section~\ref{sec:lvk}) places the fiducial Model A high
within that physically allowed band at the calibrated normalization. The BNS
rate, $55\,\Gpcyr$, is consistent both with the GWTC-5.0 BNS region and with the
COMPAS local-rate anchor to which the per-population $\langle M_{\rm evolved}\rangle$
is calibrated (Section~\ref{sec:cosmicrate}).

The engine-class assignments line up with the qualitative expectations of the
\citet{gottlieb2024} framework. That work argues that BHNS mergers should
contribute essentially only to the lbGRB plus red-kilonova population, while
sbGRBs are driven by a magnetized HMNS remnant and carry a bluer kilonova. Our
BHNS population reproduces the first point: the GRB-capable BHNS systems land
exclusively in the long-burst classes (faint lbGRB or lbGRB plus red
kilonova), never the sbGRB class. The BNS sbGRB plus blue kilonova class (24
percent) reproduces the second. Against the \citet{rastinejad2024} uniform
sample (Section~\ref{sec:obsclass}) the model recovers the observed
sbGRB-minority, long-duration-majority ordering.

Because the disk-mass classification assumes complete jet launching above
threshold, the per-class rates are upper bounds, and beaming (Eq.~\ref{eq:beam})
is what bridges them to the observed local sGRB rate. Published beaming-limited
sGRB rates span $4.1$ ($2.2$ to $6.4$)$\,\Gpcyr$ \citep{wandermanpiran2015},
$1.3$ ($0.5$ to $3.0$)$\,\Gpcyr$ \citep{ghirlanda2016}, and $3.6$ ($1.8$ to
$6.5$)$\,\Gpcyr$ \citep{colombo2022}. At the fiducial $\theta_j = 13^\circ$
the beamed sbGRB rate is $0.56\,\Gpcyr$ ($0.33$ to $0.84$ across the
$10^\circ$ to $16^\circ$ band): inside the \citet{ghirlanda2016} interval but
below the \citet{wandermanpiran2015} and \citet{colombo2022} intervals, and
the gap cannot be closed by jet-launching efficiency, which only lowers an
upper bound. The same direction of discrepancy, GW-inferred NS merger rates
that undershoot the observed sGRB rate at canonical opening angles, is the
subject of \citet{kunnumkai2026} (see also \citealt{fishbach2026,jin2026});
their proposed resolution, wider sGRB jets, acts here as well: $\theta_j
\approx 25^\circ$ raises the beamed sbGRB rate to $\approx 2\,\Gpcyr$, inside
the \citet{ghirlanda2016} and \citet{colombo2022} intervals and within 10
percent of the \citet{wandermanpiran2015} lower bound.

\subsection{Systematics across the population-synthesis grid} \label{sec:systematics}

The five-model suite of Appendix~\ref{app:robustness} isolates two assumptions; to
bound the prescription systematics more broadly we scan the full 20-model
\citet{broekgaarden2021} grid (Figures~\ref{fig:rate20} and \ref{fig:frac20}).
Two facts stand out. First, the local rate spans roughly two orders of magnitude
across the grid, from $\sim 0.6\,\Gpcyr$ for the unstable case-BB model to
$\sim 230\,\Gpcyr$ for the reduced-kick model
($\sigma_{\rm kick} = 30\,{\rm km\,s^{-1}}$), with the natal-kick,
common-envelope, and case-BB variations the dominant levers. Eighteen of the
20 BNS rates land inside the GWTC-5.0 BNS region (the unstable case-BB model
falls below it, the reduced-kick model above it). On the BHNS side the
tightened NSBH region has real discriminating power: 14 of the 20 variations
exceed it, including the entire core suite, while the $\beta = 0.5$ and
$0.75$ mass-transfer models, unstable case BB, $\alpha_{\rm CE} = 10$, and
the enhanced Wolf-Rayet-wind model ($f_{\rm WR} = 5$) land inside and the
extreme $\alpha_{\rm CE} = 0.1$ model falls below. Second, and in contrast, the GRB-class fractions are nearly
invariant across the same grid: the lbGRB-producing classes carry at least 72
percent of the BNS population in every variation, and the sbGRB-plus-blue class
stays a minority at the 11 to 28 percent level. The class demographics are therefore a robust
prediction of the framework even where the absolute rate is not, because the
classification boundaries depend on the component-mass distribution (set by the
remnant engine and the NS-mass cap) rather than on the formation rate. The
sub-structure of the grid sharpens this: the lbGRB-plus-red-kilonova (disk) class
is the one fraction that does move, growing from a few percent in the fiducial
model to roughly a quarter of the population under the $\beta = 0.5$
mass-transfer-efficiency model and to more than a third at low common-envelope
efficiency ($\alpha_{\rm CE} \leq 0.5$;
Figures~\ref{fig:alphace5} and \ref{fig:betasweep}), because both reshape the
mass-ratio distribution to which the $q > 1.2$ prompt-collapse boundary is
sensitive. The supernova remnant engine (delayed versus rapid) shifts the
component masses but leaves the class ordering intact
(Figure~\ref{fig:snengine}), and the formation-channel-to-class mapping is itself
fragile under the case-BB stability assumption (Figure~\ref{fig:chanclass}): in
the fiducial model the classic isolated-binary channel maps cleanly onto the
sbGRB-plus-blue class, but unstable case-BB mass transfer redistributes that
mapping across classes.

The equation of state is the largest single systematic on the class fractions.
The EOS$\times$model matrix (Figure~\ref{fig:eosmatrix}) shows that reading along
a row (fixed EOS, varying binary physics) barely moves the HMNS-class fraction,
whereas reading down a column (fixed binary physics, varying EOS) changes it by
factors of several through the linked $(\Mtov, \Mthresh, M_{\rm HMNS})$ boundaries.
Dense-matter uncertainty therefore dominates population-synthesis uncertainty for
the class demographics, which is the opposite of the situation for the absolute
rate.

\subsection{Sensitivity to the classification thresholds} \label{sec:thresholds}

Three modeling choices set the class boundaries and deserve explicit statement.
First, the prompt-collapse ratio $k = 1.27$ is a \citet{gottlieb2023} fiducial
chosen so that $\Mthresh \approx 2.8\,\Msun$; it sits below the
\citet{bauswein2013} EOS-derived band $k \in [1.30, 1.70]$, so our $\Mthresh$ is
conservative and a larger $k$ would move systems from the disk and faint classes
toward the HMNS class. Second, the long-lived-HMNS split at $M_{\rm HMNS} =
1.2\,\Mtov$ is a code heuristic \citep{margalit2017}, and per-event assignments
near it are EOS-sensitive (GW170817; Section~\ref{sec:massplane}). Third, the
neutron-star mass remap onto the \citet{alsing2018} double Gaussian
\citep{mandel2020} sets the absolute placement of the low-mass classes; a
different remap would move the boundary between the low- and high-$\Mtot$ classes
(Section~\ref{sec:massdist}). Because the population fractions are far more stable
than any single per-event assignment (Section~\ref{sec:systematics}), these
threshold choices change the headline fractions at the few-percentage-point
level while leaving the class ordering intact, but they dominate the
classification of individual borderline events such as GW170817.

%% ======================================================================
\FloatBarrier
\section{Conclusions} \label{sec:conclusions}
%% ======================================================================

Mapping the COMPAS compact-binary populations onto the Gottlieb engine classes
lets us ask, at the population level, which central engine each merger drives and
whether the implied rates are compatible with the gravitational-wave sample. Our
main findings are:

\begin{enumerate}
\item Applying the \citet{gottlieb2023,gottlieb2024} classification to \compas\
Model A, the BNS population divides into 24.0 percent sbGRB plus blue kilonova,
26.8 percent lbGRB plus red kilonova (HMNS), 2.2 percent lbGRB plus red kilonova
(disk), and 47.0 percent faint lbGRB; the BHNS population is 75.6 percent
NS-swallowing with no GRB.
\item The intrinsic BNS local rate ($55\,\Gpcyr$) is consistent with the
GWTC-5.0 BNS 90 percent credible region, while the BHNS rate ($84.5\,\Gpcyr$)
exceeds the tightened NSBH 90 percent region by a factor $\sim 2.6$, as does
every member of the five-model suite: the \compas\ BHNS population
generically overpredicts the GW-measured NSBH rate, most sharply at the
highest $M_{\rm NS,max}$. Within the full grid, only fixed mass-transfer
efficiency, unstable case-BB mass transfer, $\alpha_{\rm CE} = 10$, or
enhanced Wolf-Rayet winds bring the BHNS rate back inside the region.
\item Confronted with the kilonova ejecta of the \citet{rastinejad2024} uniform
sample, the long-duration bursts with kilonovae are consistent with a
black-hole accretion-disk engine ($+0.10$ dex from the null), while the
short-duration bursts lie $+3.4$ dex above it and far closer to the HMNS
disk-wind expectation; this inverts the naive duration-to-engine mapping in
the direction predicted by \citet{gottlieb2024}.
\item Across the full 20-model \citet{broekgaarden2021} grid the local rate
spans two orders of magnitude while the GRB-class fractions stay nearly
invariant, so the class demographics are a robust prediction of the framework
even where the absolute rate is not. The class ordering is in particular robust
to the common-envelope efficiency and the maximum NS mass
(Appendix~\ref{app:robustness}); the equation of state, not the binary-evolution
prescription, is the dominant systematic on the class fractions, and BH spin is
the dominant within-model source of BHNS rate uncertainty.
\end{enumerate}

%% ======================================================================
\section*{Data Availability}
%% ======================================================================

The \compas\ population-synthesis catalogs underlying this work are public:
the binary neutron star models on Zenodo (record 5189849, the
\citealt{broekgaarden2022} paper II data release) and the black hole neutron
star models on Zenodo (record 5178777, \citealt{broekgaarden2021}).
The metallicity-specific star-formation history is the IllustrisTNG TNG100-1 fit
of \citet{levina2026}. The derived quantities and figures follow from the
equations and parameters stated in Section~\ref{sec:methods}.

\begin{acknowledgments}
This work uses the \compas\ rapid binary population-synthesis catalogs of
\citet{broekgaarden2021,broekgaarden2022} (Zenodo records 5189849 and 5178777)
and the IllustrisTNG-calibrated metallicity-specific star-formation history of
\citet{levina2026}. We thank the \compas\ and IllustrisTNG collaborations for
making their data products public. The cosmic-integration machinery builds on
the \compas\ \fci\ module. [Funding, computing-allocation, and individual
acknowledgments to be completed by the authors before submission.]
\end{acknowledgments}

\software{\compas\ \citep{broekgaarden2021,broekgaarden2022},
Astropy \citep{2022ApJ...935..167A, 2018AJ....156..123A},
NumPy, SciPy, Matplotlib, pandas, h5py.}

%% ======================================================================
\appendix

\section{Formation-channel decomposition of the class rates} \label{app:channels}
%% ======================================================================

\begin{figure*}[tbp]
\epsscale{1.0}
\plotone{rate_bns_by_class_per_channel.pdf}
\caption{BNS intrinsic merger rate $\Rz$ per \citet{gottlieb2024} GRB class (2x2),
decomposed onto the \citet{broekgaarden2021} formation channels I to V; dotted
curve is the per-class total. \compas\ Model A; cosmology Planck 2015
\citep{planck2016}; rates intrinsic upper bounds. Per-curve weighted-Poisson
error is sub-percent and the common MSSFR-resolution systematic cancels in the
per-channel decomposition, so no band is drawn. \label{fig:ratechanbns}}
\end{figure*}

\begin{figure*}[tbp]
\epsscale{1.15}
\plotone{rate_bhns_by_class_per_channel.pdf}
\caption{BHNS intrinsic merger rate $\Rz$ per hybrid GRB class (1x3) at
$\aBH = 0.5$, decomposed onto the \citet{broekgaarden2021} formation channels;
dotted curve is the misalignment-corrected per-class total. \compas\ Model A;
rates intrinsic upper bounds. Curves are at $\aBH = 0.5$; the BH-spin systematic
(Figure~\ref{fig:spinsweep}) is the dominant per-class uncertainty and is
not propagated per channel. \label{fig:ratechanbhns}}
\end{figure*}

Figures~\ref{fig:ratechanbns} and \ref{fig:ratechanbhns} decompose the
per-class rates of Section~\ref{sec:classrates} onto the
\citet{broekgaarden2021} formation channels of Section~\ref{sec:channels}. The decomposition
satisfies two additivity constraints to a relative tolerance of $10^{-12}$: summing
over channels recovers the per-class rate, and summing over classes within a
channel recovers the per-channel rate. For BNS the rate is carried by channels~I
(stable mass transfer plus common envelope) and~IV (double-core common
envelope) across all classes; channel~IV dominates the faint class because its
tighter post-common-envelope separations drive the shortest delay times. For
BHNS the rate is almost entirely channel~I, with the misalignment factor
applied per cell.

\FloatBarrier
%% ======================================================================
\section{Population-synthesis variations} \label{app:robustness}
%% ======================================================================

To test how strongly the class predictions depend on the binary-evolution
assumptions, we repeat the BNS classification across the five
\citet{broekgaarden2021} models, recalibrating $M_{\rm evolved}$ per model and
population. The variations isolate two parameters: the common-envelope efficiency
($\alpha_{\rm CE}$: Models F $=0.5$, A $=1.0$, G $=2.0$) and the maximum
neutron-star mass ($M_{\rm NS,max}$: Models J $=2.0$, A $=2.5$, K $=3.0\,\Msun$).
The natal-kick variations (Models M, N, O) enter at the end of this appendix,
where they test the host-galaxy offset prediction.

The resulting BNS class fractions appear per model in the J, A, and K rows of
Figure~\ref{fig:frac20} for the maximum NS mass, and as the five-point
$\alpha_{\rm CE}$ sequence of Figure~\ref{fig:alphace5} (which contains Models
F, A, and G) for the common-envelope efficiency. The lbGRB plus red kilonova
(HMNS) class is present across all variations, and the fractions shift smoothly
with both parameters, so the qualitative class ordering of
Section~\ref{sec:massdist} is robust to these assumptions. Model J has
$M_{\rm NS,max} = 2.0\,\Msun < \Mtov = 2.2\,\Msun$, so its $[2.0, 2.2]\,\Msun$
window is \compas-BH but EOS-NS; we retain the \compas\ labels throughout for
consistency.

Figure~\ref{fig:chanfrac} shows the per-model formation-channel fractions. The
common-envelope efficiency rebalances channels~I and~IV for BNS, and the
low-efficiency Model F carries the largest stable-mass-transfer-only
(channel~II) BHNS share, but channel~I remains
dominant for both populations across the suite.

\begin{figure*}[tbp]
\epsscale{1.15}
\plotone{channel_fractions_per_model.pdf}
\caption{Per-model \citet{broekgaarden2021} formation-channel fractions (I to V)
for BNS (left) and BHNS (right) across Models A, F, G, J, K. Colors follow the
project model palette (matplotlib cividis); colormap choice follows
\citet{crameri2020}. \label{fig:chanfrac}}
\end{figure*}

The per-model local rates against the GWTC-5.0 90 percent credible regions
\citep{abac2026} appear in Figure~\ref{fig:rate20}, which spans the full grid
and contains the core suite. The BNS rate stays within the BNS region
$[5.1, 154.7]\,\Gpcyr$ for all five models: $55.0$ (A), $29.3$ (F), $106.0$ (G),
$53.7$ (J), and $56.9\,\Gpcyr$ (K). Every BHNS rate, in contrast, exceeds the
tightened NSBH region $[6.7, 32.8]\,\Gpcyr$: $38.0$ (F, $\sim 1.2\times$ the
upper bound), $68.9$ (J), $78.1$ (G), $84.5$ (the fiducial Model A,
$\sim 2.6\times$), and $96.3\,\Gpcyr$ (K, $\sim 2.9\times$;
$M_{\rm NS,max} = 3.0\,\Msun$). The overprediction noted in
Section~\ref{sec:lvk} therefore applies across the whole suite, deepening with
increasing maximum NS mass.

Figure~\ref{fig:delaycdfmodel} shows the weighted BNS delay-time CDFs. The
common-envelope efficiency shifts the distribution, with Model F at much longer
delays \citep{vanson2022,broekgaarden2022}, whereas the maximum-NS-mass
variations J and K leave it nearly unchanged (K acquires only a mild
long-delay excess), as expected for a
remnant-mass cap that does not alter orbital separations.

\begin{figure}[tbp]
\epsscale{1.15}
\plotone{delay_time_cdf_by_model.pdf}
\caption{Weighted BNS delay-time CDFs ($t_{\rm delay} = t_{\rm form} + t_c$, in
Myr) for \citet{broekgaarden2021} Models A, F, G, J, K. Model F
($\alpha_{\rm CE} = 0.5$) shifts to markedly longer delays because the
short-delay double-core-CE channel does not survive at low efficiency
(Figure~\ref{fig:chanfrac}); the maximum-NS-mass
variations (J, K) leave the curve nearly unchanged. \compas,
STROOPWAFEL-weighted.
\label{fig:delaycdfmodel}}
\end{figure}

The natal-kick variations test the one main-text prediction that does not
reduce to the component-mass distribution, the host-galaxy offsets.
Figure~\ref{fig:offsetkicks} repeats the offset prediction of
Section~\ref{sec:offsets} for the reduced core-collapse kick dispersions
($\sigma_{\rm kick} = 100$ and $30\,{\rm km\,s^{-1}}$, Models M and N) and,
for BHNS, the no-BH-natal-kick Model O. The offset CDFs are nearly insensitive
to the prescription: changing the kick dispersion reshapes which binaries stay
bound and merge rather than simply rescaling the survivors' systemic
velocities, so the offsets remain a tracer of the host potential and the
delay-time distribution rather than of the kick model.

\begin{figure*}[tbp]
\epsscale{1.15}
\plotone{projected_offsets_kick_models.pdf}
\caption{Weighted projected galactocentric offset CDFs (BNS left, BHNS right)
across the natal-kick variations of Table~\ref{tab:models}: the fiducial
\citet{hobbs2005} $\sigma_{\rm kick} = 265\,{\rm km\,s^{-1}}$ Maxwellian
(Model A), reduced dispersions of $100$ (Model M) and $30\,{\rm km\,s^{-1}}$
(Model N), and no BH natal kick (Model O, BHNS only). Orbit integration, host
mix, and observed comparators as in Figure~\ref{fig:offsets}. \compas,
STROOPWAFEL-weighted. \label{fig:offsetkicks}}
\end{figure*}

Figure~\ref{fig:eosmatrix} crosses the five models with the four EOS, reporting
the weighted lbGRB plus red kilonova (HMNS) fraction. Each EOS sets $\Mtov$ and
the linked $\Mthresh = k\,\Mtov$, while the population synthesis sets the BNS
mass distribution. The HMNS class is present for every (model, EOS) pair,
confirming that the engine ordering survives the joint variation of
binary-evolution and dense-matter assumptions.

\begin{figure}[tbp]
\epsscale{1.15}
\plotone{eos_x_model_matrix.pdf}
\caption{Weighted lbGRB plus red kilonova (HMNS) class fraction on the
$5 \times 4$ grid of \citet{broekgaarden2021} models (rows A, F, G, J, K) and
EOS (columns APR4, SFHo, LS220, DD2; constants from \citealt{bauswein2013,
bauswein2020} and \citealt{read2009}). Each EOS fixes $\Mtov$ and the linked
$\Mthresh = k\,\Mtov$. Color follows the perceptually uniform cividis map;
colormap choice follows \citet{crameri2020}. \label{fig:eosmatrix}}
\end{figure}

\FloatBarrier
\section{The full 20-model grid} \label{app:grid}

Beyond the five-model core suite, the complete
\citet{broekgaarden2021,broekgaarden2022} release spans 20 binary-physics
variations, collected with their assumptions in Table~\ref{tab:models}. We
stream every variation through the same classification and cosmic-integration
pipeline, caching the per-model local rates, class fractions, and
weighted-Poisson errors. Figure~\ref{fig:rate20}
gives the per-model local rates against the GWTC-5.0 regions, and
Figure~\ref{fig:frac20} the per-model class fractions; together they show the
two-orders-of-magnitude rate spread against the near-invariant class composition
discussed in Section~\ref{sec:systematics}. Figures~\ref{fig:alphace5} to
\ref{fig:chanclass} isolate the individual assumptions that do move the fractions: the
five-point common-envelope sequence, the mass-transfer-efficiency sweep, the
delayed-versus-rapid supernova engine, and the channel-to-class contingency under
case-BB and common-envelope-survival variations.

\begin{figure*}[tbp]
\epsscale{0.9}
\plotone{rate_local_20models.pdf}
\caption{Per-model local intrinsic merger rate $\Rzero$ for all 20
\citet{broekgaarden2021} variations, BNS (left) and BHNS (right), with the core
suite (A, F, G, J, K) first and the remaining variations in registry order
(Table~\ref{tab:models}), against the GWTC-5.0 90 percent credible regions
\citep[shaded]{abac2026}. The weighted-Poisson sampling error is sub-percent and
the MSSFR-resolution systematic ($\sim \pm 30$ percent, common to all bars) is
smaller than the cross-model spread; all GRB rates are intrinsic upper bounds.
\label{fig:rate20}}
\end{figure*}

\begin{figure*}[tbp]
\epsscale{0.9}
\plotone{class_fractions_20models.pdf}
\caption{Weighted \citet{gottlieb2024} class fractions and local rates for
all 20 \citet{broekgaarden2021} variations, BNS four-class (left block) and
BHNS three-class (right block), following the presentation of
\citet{broekgaarden2026}. Within each block the stacked bars give the
weighted class fractions, with models sorted by the first-stacked class
fraction (sbGRB plus blue kilonova for BNS, lbGRB plus red kilonova for
BHNS), largest on top, rather than in registry order; the narrow companion
panel gives the same models' local intrinsic rate $\Rzero$ (diamonds, log
axis) against the GWTC-5.0 90 percent credible region
\citep[shaded]{abac2026}, with rates as in Figure~\ref{fig:rate20}. The
weighted-Poisson sampling error on each fraction is sub-percent; the
composition is far more stable across the grid than the rate.
\label{fig:frac20}}
\end{figure*}

\begin{figure*}[tbp]
\epsscale{1.15}
\plotone{class_fractions_alpha_ce_5pt.pdf}
\caption{Five-point common-envelope sequence $\alpha_{\rm CE} \in \{0.1, 0.5,
1.0, 2.0, 10\}$. \textit{Left:} stacked BNS class fractions.
\textit{Right:} BNS and BHNS local rate versus $\alpha_{\rm CE}$; the
rate points share the MSSFR-resolution systematic of Figure~\ref{fig:rate20}
($\sim \pm 30$ percent, common to all points and smaller than the trends shown,
so no whisker is drawn).
\label{fig:alphace5}}
\end{figure*}

\begin{figure}[tbp]
\epsscale{1.15}
\plotone{class_fractions_beta_sweep.pdf}
\caption{Stacked BNS class fractions for the fiducial adaptive mass-transfer
efficiency (Model A) and fixed $\beta \in \{0.25, 0.5, 0.75\}$ (Models B, C, D).
The lbGRB-plus-red-kilonova (disk) class grows to roughly a quarter of the
population at $\beta = 0.5$. \label{fig:betasweep}}
\end{figure}

\begin{figure*}[tbp]
\epsscale{1.15}
\plotone{class_fractions_sn_engine.pdf}
\caption{Delayed (Model A) versus rapid (Model I) supernova remnant engine.
\textit{Left:} BNS class fractions. \textit{Center:}
weighted BNS total-mass distribution. \textit{Right:} weighted BHNS BH-mass
distribution. The rapid engine enforces the lower mass gap but leaves the class
ordering intact. \label{fig:snengine}}
\end{figure*}

\begin{figure*}[tbp]
\epsscale{0.85}
\plotone{channel_class_casebb_ce.pdf}
\caption{Weighted $P(\mathrm{class}\,|\,\mathrm{channel})$ for the fiducial Model
A and three variations that reshape the formation-channel mix: E (unstable case
BB), EH (case BB plus optimistic common-envelope survival), and H (optimistic
common-envelope survival). Rows are \citet{broekgaarden2021} channels I to V,
columns are \citet{gottlieb2024} BNS classes; rows sum to one. The
channel-to-class mapping that is clean in the fiducial model is redistributed
under unstable case-BB mass transfer. \label{fig:chanclass}}
\end{figure*}

\clearpage

\bibliographystyle{aasjournalv7.1}
\bibliography{grb_paper}

\end{document}
